\section{The nonabelian Hodge homeomorphism}

In \cite[\S 8,9]{simpson-local_systems_on_proper} Simpson showed, using a simplicial resolution, that the moduli spaces of vector bundles over a smooth Deligne-Mumford stack respect the Riemann-Hilbert correspondence and the non-abelian Hodge homeomorphism.

More specifically, for a smooth Deligne-Mumford stack $\mc{Y}$ he established the Riemann-Hilbert isomorphism
\begin{equation} \label{riemann hilbert DM stack}
M_{n, DR}(\mc{Y})^{\text{an}} \cong M_{n, B}(\mc{Y})^{\text{an}},
\end{equation}
and the non-abelian Hodge homeomorphism
\begin{equation}\label{non abelian Hodge DM stack}
M_{n, Dol}(\mc{Y})^{\text{top}} \cong M_{n, B}(\mc{Y})^{\text{top}}.
\end{equation}

Let $\mc{Y} = \mc{X}$ be the stacky curve that we have studied so far. Of course, \eqref{riemann hilbert DM stack} and \eqref{non abelian Hodge DM stack} respect the decomposition in connected components, so they restrict to
\begin{gather}
M_{\underline{m}, DR}(\mc{X})^{\text{an}} \cong M_{\underline{m}, B}(\mc{Y})^{\text{an}}, \\
M_{\underline{m}, Dol}^0(\mc{X})^{\text{top}} \cong M_{\underline{m}, B}(\mc{X})^{\text{top}}.
\end{gather}

Using the observations made in \zcref{section: de rham moduli space}, these translate to homeomorphisms
\begin{equation}
M_{\underline{m}, Dol}^0(\mc{X})^{\text{top}} \cong \left( \bigg\{ (A_i, B_i)_{i = 1}^g \in GL_n^{2g +1} \;\bigg|\; \prod_{i=1}^{g}[A_i, B_i] = C_{\underline{m}} \bigg\} // GL_n \right)^{\text{top}},
\end{equation}
and in particular to
\begin{equation}
M_{n_{-d}, Dol}^0(\mc{X})^{\text{top}} \cong \left( \bigg\{ (A_i, B_i)_{i = 1}^g \in GL_n^{2g +1} \,\bigg|\, \prod_{i=1}^{g}[A_i, B_i] = \mathrm{e}^{-\frac{2 \pi \mathrm{i}}{n} d} Id \bigg\} // GL_n \right)^{\text{top}}.
\end{equation}

On the other hand, the equivalence \eqref{eq: equivalence of dolbeault moduli spaces} proven in \zcref{section: dolbeault moduli spaces} induces an isomorphism of the coarse moduli spaces
\begin{equation}
M_{n, Dol}^d(X) \cong M_{n_{-d}, Dol}^0(\mc{X}).
\end{equation}

Finally, putting these together we get an homeomorphism
\begin{equation}
M_{n, Dol}^d(X)^{\text{top}} \cong \left( \bigg\{ (A_i, B_i)_{i = 1}^g \in GL_n^{2g +1} \,\bigg|\, \prod_{i=1}^{g}[A_i, B_i] = \mathrm{e}^{-\frac{2 \pi \mathrm{i}}{n} d} Id \bigg\} // GL_n \right)^{\text{top}},
\end{equation}
therefore the space on the right hand side can be interpreted as version of the character variety for bundles of nonzero degree on $X$.

The picture can be explained with more clarity by embedding both $M_{n, Dol}^d$ and $M_{n, DR}^d$ in a bigger moduli space $M_{n, Hdg}^d$. Indeed, for a fixed complex number $\lambda$, define the notion of $\lambda$-connection on a vector bundle $\mc{E}$ as a morphism $\nabla\colon \mc{E} \to \mc{E} \otimes \Omega$, respecting the Leibniz rule up to the scalar $\lambda$, namely
\[
\nabla(f \cdot s) = s \otimes \dd f + \lambda f \cdot \nabla(s).
\]
Then $M_{n, Hdg}^d(\mc{X})$ is the moduli space parametrising isomorphism classes of triples $(\mc{E}, \nabla, \lambda)$. Notice that
\begin{itemize}
	\item the multplicative group $\Gb_m$ acts on $M_{n, Hdg}^{d}$ by
	\[
	t \cdot (\Ec, \nabla) = (\Ec, t \nabla);
	\]
	\item there is a $\Gb_m$-equivariant projection
	\begin{align}
	\pi\colon M_{n, Hdg}^d &\to \Ab^1 \\
	(\mc{E}, \nabla, \lambda) &\mapsto \lambda,
	\end{align}
	where $\Gb_m$ acts on $\Ab^1$ with weight one;
	\item the fibers $\pi^{-1}(\lambda)$ are all isomorphic for $\lambda \neq 0$ to $\pi^{-1}(1) = M_{n, DR}^d$, while $pi^{-1}(0)= M_{n, Dol}^d$.
\end{itemize}
This picture restricts to any connected component, determined by numerical invariants. Summing up, letting $M_{Hdg} \coloneqq M_{n_{-d}, Hdg}^0(\mc{X}), M_{Dol} \coloneqq M_{n_{-d}, Dol}^0(\mc{X})$ and $M_{Dr} \coloneqq M_{n_{-d}, DR}^0(\mc{X})$, we have the following diagram
\begin{equation}
\begin{tikzcd}
	{M_{Dol}} & {M_{Hdg}} & {M_{DR}} & {M_B} \\
	{\{ 0 \}} & {\Ab^1} & {\{ 1 \},}
	\arrow["i", hook, from=1-1, to=1-2]
	\arrow[from=1-1, to=2-1]
	\arrow["\lrcorner"{anchor=center, pos=0.125}, draw=none, from=1-1, to=2-2]
	\arrow["\pi", from=1-2, to=2-2]
	\arrow["j"', hook', from=1-3, to=1-2]
	\arrow["\psi", "\cong"', from=1-3, to=1-4]
	\arrow["\lrcorner"{anchor=center, pos=0.125, rotate=-90}, draw=none, from=1-3, to=2-2]
	\arrow[from=1-3, to=2-3]
	\arrow[from=2-1, to=2-2]
	\arrow[from=2-3, to=2-2]
\end{tikzcd}
\end{equation}
where the inclusions $i$ and $j$ are algebraic, while $\psi$ is an isomorphism of complex manifold that is only defined analytically via the Riemann-Hilbert correspondence \eqref{riemann hilbert DM stack}.

We want to compare the singular cohomology with rational coefficients of $M_{Dol}, \; M_{Hdg}, \; M_{DR}$. The idea to use the Bia\l{y}nicki-Birula decomposition was inspired by Nakajima's appendix to \cite{appendix-by-nakajima}.
%Nakajima used this argument for point counting; since we want to compute cohomology instead, we will need some extra steps.

Notice that $M_{Hdg}$ is a \textit{semiprojective variety}, according to the terminology used in \cite{Hausel-RodriguezVillega}. This means that it is a quasi-projective variety with a $\Gb_m$-action such that
\begin{itemize}
    \item the fixed point locus $M_{Hdg}^{\Gb_m}$ is proper;
    \item the limit point $\displaystyle \lim_{t \to 0} t \cdot x$ exists for all $x \in M_{Hdg}$.
\end{itemize}
Indeed, these properties are \cite[Lemma 10.4, Corollary 10.3]{Simpson-Hodge-filtration}.

Then the picture above respects the assumptions of \cite[Corollary 1.3.3]{Hausel-RodriguezVillega}, proving the following theorem.

\begin{thm}
The inclusions $i$ and $j$ induce isomorphisms of mixed hodge structures in cohomology
\[
H^k(M_{Dol}) \cong H^k(M_{Hdg}) \cong H^k(M_{DR}).
\]
\end{thm}

Now we are left to try and understand how $\psi$ acts on the mixed Hodge structures. A first step in this direction, is to choose generators of pure weights for the various cohomology rings and track their weight along the various isomorphisms. In order to do this, we take the universal bundle $(\mc{E}_\text{univ}, \nabla_\text{univ})$ on $M_{Hdg}$.

\subsection{Comparing the cohomologies of the moduli spaces}

If we want to compare the cohomology of the moduli spaces, it is more useful to try and use algebraic maps, instead of the homeomorphism \zcref{non abelian Hodge DM stack}, since the use of algebraic maps allows to compare Hodge structures.

Let $M \coloneqq M_{n_{-d}, Hdg}^0(\mc{X}), M_0 \coloneqq M_{n_{-d}, Dol}^0(\mc{X})$ and $M_1 \coloneqq M_{n_{-d}, DR}^0(\mc{X})$.
We have the following pullback diagram
\[\begin{tikzcd}
	{M_0} & M & {M_{\neq 0} \coloneqq M \setminus M_0} \\
	{\{ 0 \}} & {\Ab^1} & {\Gb_m.}
	\arrow["i", hook, from=1-1, to=1-2]
	\arrow[from=1-1, to=2-1]
	\arrow["\lrcorner"{anchor=center, pos=0.125}, draw=none, from=1-1, to=2-2]
	\arrow["\pi", from=1-2, to=2-2]
	\arrow["j"', hook', from=1-3, to=1-2]
	\arrow["\lrcorner"{anchor=center, pos=0.125, rotate=-90}, draw=none, from=1-3, to=2-2]
	\arrow[from=1-3, to=2-3]
	\arrow[from=2-1, to=2-2]
	\arrow[from=2-2, to=2-3]
\end{tikzcd}\]
Moreover, $M$ is endowed with a $\Gb_m$-action induced by
\[
t \cdot (\Ec, \nabla) = (\Ec, t \nabla),
\]
making $\pi$ into a $\Gb_m$-equivariant morphism. %[easy computation, may add]

We want to compare the singular cohomology with rational coefficients of the complex manifolds obtained from $M, M_0, M_1, M_{\neq 0}$. The idea to use the Bia\l{y}nicki-Birula decomposition was inspired by Nakajima's appendix to \cite{appendix-by-nakajima}. Nakajima used this argument for point counting; since we want to compute cohomology instead, we will need some extra steps.

Notice that $M$ is a \textit{semiprojective variety}, according to the terminology used in \cite{Hausel-RodriguezVillega}. This means that it is a quasi-projective variety with a $\Gb_m$-action such that:
\begin{itemize}
    \item the fixed point locus $M^{\Gb_m}$ is proper;
    \item the limit point $\displaystyle \lim_{t \to 0} t \cdot x$ exists for all $x \in M$.
\end{itemize}
Indeed, these properties are \cite[Lemma 10.4, Corollary 10.3]{Simpson-Hodge-filtration}.

Then the picture above respects the assumptions of \cite[Corollary 1.3.3]{Hausel-RodriguezVillega}, proving the following theorem.

\begin{thm}
The inclusions $i$ and $j$ induce isomorphisms of mixed hodge structures in cohomology
\[
H^k(M_0) \cong H^k(M) \cong H^k(M_1).
\]
\end{thm}